\documentclass[11pt,a4paper]{article}
\usepackage[utf8]{inputenc}
\usepackage[T1]{fontenc}
\usepackage[catalan]{babel}
\usepackage{lmodern}
\usepackage[sfdefault,lining]{FiraSans}
\usepackage[a4paper,top=2.3cm,bottom=2.2cm,left=1.8cm,right=1.8cm,headheight=1.3cm,headsep=0.4cm,footskip=1cm]{geometry}
\usepackage{amsmath,amssymb}
\usepackage{icomma}
\usepackage{graphicx}
\usepackage[table]{xcolor}
\usepackage{tikz}
\usepackage{fontawesome5}
\usepackage{ccicons}
\usepackage{fancyhdr}
\usepackage{lastpage}
\usepackage{tasks}
\usepackage{needspace}
\usepackage[most]{tcolorbox}
\usepackage{array}
\usepackage[hidelinks,pdftitle={Fitxa d'exercicis de la unitat 11: Cossos geomètrics. 3r d'ESO},pdfsubject={Matemàtiques 3r ESO},pdfkeywords={poliedres, fórmula d'Euler, poliedres regulars, prisma, piràmide, cilindre, con, esfera, àrea, volum, capacitat}]{hyperref}

\definecolor{navy}{HTML}{1B2748}
\definecolor{accent}{HTML}{E8603A}
\definecolor{teal}{HTML}{1F9E90}
\definecolor{soft}{HTML}{EEF1F8}
\definecolor{muted}{HTML}{5D6685}

\setlength{\parindent}{0pt}
\setlength{\parskip}{3pt}
\settasks{label=\alph*), label-format=\bfseries\color{navy}, label-width=1.3em, item-indent=1.8em, column-sep=0.6em, after-item-skip=5pt, before-skip=3pt, after-skip=4pt}

% ---------- capçalera i peu ----------
\newcommand{\llicencia}{\href{https://creativecommons.org/licenses/by-nc-sa/4.0/deed.ca}{\ccbyncsa}\ \ Llicència Creative Commons Reconeixement-NoComercial-CompartirIgual 4.0 Internacional (CC BY-NC-SA 4.0)}
\pagestyle{fancy}
\fancyhf{}
\fancyhead[L]{\raisebox{-0.35\height}{\includegraphics[height=0.95cm]{logo-gva.png}}}
\fancyhead[C]{\small\color{muted}Fitxa d'exercicis · Unitat 11: Cossos geomètrics · 3r d'ESO}
\fancyhead[R]{\raisebox{-0.35\height}{\includegraphics[height=0.95cm]{logo-vermellar.png}}}
\fancyfoot[L]{\scriptsize\color{muted}\llicencia}
\fancyfoot[R]{\small\color{muted}\thepage}
\renewcommand{\headrulewidth}{0.4pt}
\renewcommand{\headrule}{\color{muted!40}\hrule height \headrulewidth}
\fancypagestyle{primera}{\fancyhf{}\renewcommand{\headrulewidth}{0pt}
  \fancyfoot[L]{\scriptsize\color{muted}\llicencia}
  \fancyfoot[R]{\small\color{muted}\thepage}}

% ---------- nivells i exercicis ----------
\newcommand{\estrelles}[1]{{\color{accent}\ifcase#1\or\faStar\faStar[regular]\faStar[regular]\or\faStar\faStar\faStar[regular]\or\faStar\faStar\faStar\fi}}
\newcounter{exer}
\newcommand{\ex}[1]{\par\needspace{4\baselineskip}\medskip\stepcounter{exer}\noindent{\color{navy}\bfseries\theexer.}\ \ {\small\estrelles{#1}}\ \ }

\newcommand{\seccio}[2]{\par\needspace{10\baselineskip}\bigskip
  \begin{tcolorbox}[colback=navy,colframe=navy,coltext=white,arc=2mm,boxrule=0pt,left=3mm,right=3mm,top=1.2mm,bottom=1.2mm]
  {\large\bfseries #1\quad #2}
  \end{tcolorbox}\vspace{-2pt}}

\newcolumntype{C}[1]{>{\centering\arraybackslash}p{#1}}
\usepackage{pifont}

% quadre de recordatori (només a la fitxa d'exercicis)
\newcommand{\recorda}[1]{\par\smallskip\begin{tcolorbox}[colback=accent!7,colframe=accent!55,arc=2mm,boxrule=0.5pt,left=2.5mm,right=2.5mm,top=1mm,bottom=1mm]\small{\color{accent}\faLightbulb[regular]}\ \textbf{Recorda:}\ #1\end{tcolorbox}}
\newcommand{\res}[1]{#1}
\newcommand{\rest}[1]{#1}

\def\fmneed{14}
% =====================================================================
% Capçalera FentMates de les fitxes: un color per a cada curs.
% Els colors són els de la pàgina inicial (LEVELS a assets/js/materials.js);
% el logo ve de _eines/logo/pdf/ (python3 _eines/logo/genera.py).
%
% Ús, al final del preàmbul de la unitat:
%   % =====================================================================
% Capçalera FentMates de les fitxes: un color per a cada curs.
% Els colors són els de la pàgina inicial (LEVELS a assets/js/materials.js);
% el logo ve de _eines/logo/pdf/ (python3 _eines/logo/genera.py).
%
% Ús, al final del preàmbul de la unitat:
%   % =====================================================================
% Capçalera FentMates de les fitxes: un color per a cada curs.
% Els colors són els de la pàgina inicial (LEVELS a assets/js/materials.js);
% el logo ve de _eines/logo/pdf/ (python3 _eines/logo/genera.py).
%
% Ús, al final del preàmbul de la unitat:
%   \input{../comu/fentmates.tex}
%   \fmcurs{1eso}
% i, al començament del document:
%   \capcalera{Matemàtiques}{1r ESO}{Fitxa d'exercicis · Unitat 2}{Divisibilitat}
% Opcional abans de \input: \def\fmneed{14} (espai mínim abans de cada secció).
% =====================================================================
\makeatletter
\def\fm@logodir{../../logo/pdf}
\providecommand{\fmneed}{10}
\def\fm@cursnom{}

% clau del curs, color i color del text damunt del color
\newcommand{\fm@curs}[3]{\@namedef{fm@color@#1}{#2}\@namedef{fm@text@#1}{#3}}
\fm@curs{1eso}{E6B84A}{navy}
\fm@curs{2eso}{3B82F6}{white}
\fm@curs{3eso}{FF7A50}{white}
\fm@curs{4eso-a}{A78BFA}{white}
\fm@curs{4eso-b}{6D28D9}{white}
\fm@curs{1batx-mi}{9B1C1C}{white}
\fm@curs{1batx-ccss}{EF6B6B}{white}
\fm@curs{2batx-mii}{1E6B47}{white}
\fm@curs{2batx-ccss}{2BB3A3}{white}

\newcommand{\fmcurs}[1]{%
  \@ifundefined{fm@color@#1}{\PackageError{fentmates}{Curs desconegut: #1}{}}{}%
  \edef\fm@tmp{\noexpand\definecolor{curs}{HTML}{\@nameuse{fm@color@#1}}%
               \noexpand\colorlet{curstext}{\@nameuse{fm@text@#1}}}\fm@tmp
  \colorlet{cursfosc}{curs!72!black}%
  \def\fm@logo{\fm@logodir/vertical-#1.pdf}%
  \def\fm@marca{\fm@logodir/marca-#1.pdf}}

% ---------- capçalera de la primera pàgina ----------
% #1 matèria (petit), #2 curs, #3 tipus de document i unitat, #4 títol
\newcommand{\capcalera}[4]{%
  \gdef\fm@cursnom{#2}%
  \par\noindent
  \begin{tcolorbox}[colback=curs,colframe=curs,boxrule=0pt,arc=2.2mm,width=3.5cm,height=2.15cm,
      valign=center,left=3mm,right=2mm,top=1mm,bottom=1mm,nobeforeafter,box align=center]
    \color{curstext}{\footnotesize #1}\par\vspace{1pt}{\fontsize{19}{21}\selectfont\bfseries #2\par}%
  \end{tcolorbox}\hspace{5mm}%
  \begin{minipage}[c]{\dimexpr\linewidth-3.5cm-5mm-2.9cm\relax}
    {\color{cursfosc}\small\bfseries\MakeUppercase{#3}}\par\vspace{2pt}
    {\color{navy}\fontsize{20}{23}\selectfont\bfseries #4\par}
  \end{minipage}\hfill
  \begin{minipage}[c]{2.6cm}\raggedleft\includegraphics[width=2.6cm]{\fm@logo}\end{minipage}\par
  \vspace{2pt}}

% ---------- barres de secció del color del curs ----------
\renewcommand{\seccio}[2]{\par\needspace{\fmneed\baselineskip}\bigskip
  \begin{tcolorbox}[colback=curs,colframe=curs,coltext=curstext,arc=2mm,boxrule=0pt,left=3mm,right=3mm,top=1.2mm,bottom=1.2mm]
  {\large\bfseries #1\quad #2}
  \end{tcolorbox}\vspace{-2pt}}

% ---------- capçalera i peu de les pàgines ----------
% primera pàgina: logos institucionals al peu
\fancypagestyle{primera}{\fancyhf{}\renewcommand{\headrulewidth}{0pt}
  \fancyfoot[L]{\raisebox{-0.32\height}{\includegraphics[height=0.78cm]{logo-gva.png}}\hspace{5mm}%
    \raisebox{-0.32\height}{\includegraphics[height=0.78cm]{logo-vermellar.png}}}
  \fancyfoot[R]{\scriptsize\color{muted}\href{https://creativecommons.org/licenses/by-nc-sa/4.0/deed.ca}{\ccbyncsa}\ \ CC BY-NC-SA 4.0 · fentmates.org\qquad\small\thepage}}
% la resta: el símbol a l'esquerra i el curs a la dreta, amb el títol al centre com abans
\fancyhead[L]{\raisebox{-0.3\height}{\includegraphics[height=0.62cm]{\fm@marca}}}
\fancyhead[R]{\small\bfseries\color{cursfosc}\fm@cursnom}
\renewcommand{\headrule}{\color{curs!45}\hrule height \headrulewidth}
\makeatother

%   \fmcurs{1eso}
% i, al començament del document:
%   \capcalera{Matemàtiques}{1r ESO}{Fitxa d'exercicis · Unitat 2}{Divisibilitat}
% Opcional abans de \input: \def\fmneed{14} (espai mínim abans de cada secció).
% =====================================================================
\makeatletter
\def\fm@logodir{../../logo/pdf}
\providecommand{\fmneed}{10}
\def\fm@cursnom{}

% clau del curs, color i color del text damunt del color
\newcommand{\fm@curs}[3]{\@namedef{fm@color@#1}{#2}\@namedef{fm@text@#1}{#3}}
\fm@curs{1eso}{E6B84A}{navy}
\fm@curs{2eso}{3B82F6}{white}
\fm@curs{3eso}{FF7A50}{white}
\fm@curs{4eso-a}{A78BFA}{white}
\fm@curs{4eso-b}{6D28D9}{white}
\fm@curs{1batx-mi}{9B1C1C}{white}
\fm@curs{1batx-ccss}{EF6B6B}{white}
\fm@curs{2batx-mii}{1E6B47}{white}
\fm@curs{2batx-ccss}{2BB3A3}{white}

\newcommand{\fmcurs}[1]{%
  \@ifundefined{fm@color@#1}{\PackageError{fentmates}{Curs desconegut: #1}{}}{}%
  \edef\fm@tmp{\noexpand\definecolor{curs}{HTML}{\@nameuse{fm@color@#1}}%
               \noexpand\colorlet{curstext}{\@nameuse{fm@text@#1}}}\fm@tmp
  \colorlet{cursfosc}{curs!72!black}%
  \def\fm@logo{\fm@logodir/vertical-#1.pdf}%
  \def\fm@marca{\fm@logodir/marca-#1.pdf}}

% ---------- capçalera de la primera pàgina ----------
% #1 matèria (petit), #2 curs, #3 tipus de document i unitat, #4 títol
\newcommand{\capcalera}[4]{%
  \gdef\fm@cursnom{#2}%
  \par\noindent
  \begin{tcolorbox}[colback=curs,colframe=curs,boxrule=0pt,arc=2.2mm,width=3.5cm,height=2.15cm,
      valign=center,left=3mm,right=2mm,top=1mm,bottom=1mm,nobeforeafter,box align=center]
    \color{curstext}{\footnotesize #1}\par\vspace{1pt}{\fontsize{19}{21}\selectfont\bfseries #2\par}%
  \end{tcolorbox}\hspace{5mm}%
  \begin{minipage}[c]{\dimexpr\linewidth-3.5cm-5mm-2.9cm\relax}
    {\color{cursfosc}\small\bfseries\MakeUppercase{#3}}\par\vspace{2pt}
    {\color{navy}\fontsize{20}{23}\selectfont\bfseries #4\par}
  \end{minipage}\hfill
  \begin{minipage}[c]{2.6cm}\raggedleft\includegraphics[width=2.6cm]{\fm@logo}\end{minipage}\par
  \vspace{2pt}}

% ---------- barres de secció del color del curs ----------
\renewcommand{\seccio}[2]{\par\needspace{\fmneed\baselineskip}\bigskip
  \begin{tcolorbox}[colback=curs,colframe=curs,coltext=curstext,arc=2mm,boxrule=0pt,left=3mm,right=3mm,top=1.2mm,bottom=1.2mm]
  {\large\bfseries #1\quad #2}
  \end{tcolorbox}\vspace{-2pt}}

% ---------- capçalera i peu de les pàgines ----------
% primera pàgina: logos institucionals al peu
\fancypagestyle{primera}{\fancyhf{}\renewcommand{\headrulewidth}{0pt}
  \fancyfoot[L]{\raisebox{-0.32\height}{\includegraphics[height=0.78cm]{logo-gva.png}}\hspace{5mm}%
    \raisebox{-0.32\height}{\includegraphics[height=0.78cm]{logo-vermellar.png}}}
  \fancyfoot[R]{\scriptsize\color{muted}\href{https://creativecommons.org/licenses/by-nc-sa/4.0/deed.ca}{\ccbyncsa}\ \ CC BY-NC-SA 4.0 · fentmates.org\qquad\small\thepage}}
% la resta: el símbol a l'esquerra i el curs a la dreta, amb el títol al centre com abans
\fancyhead[L]{\raisebox{-0.3\height}{\includegraphics[height=0.62cm]{\fm@marca}}}
\fancyhead[R]{\small\bfseries\color{cursfosc}\fm@cursnom}
\renewcommand{\headrule}{\color{curs!45}\hrule height \headrulewidth}
\makeatother

%   \fmcurs{1eso}
% i, al començament del document:
%   \capcalera{Matemàtiques}{1r ESO}{Fitxa d'exercicis · Unitat 2}{Divisibilitat}
% Opcional abans de \input: \def\fmneed{14} (espai mínim abans de cada secció).
% =====================================================================
\makeatletter
\def\fm@logodir{../../logo/pdf}
\providecommand{\fmneed}{10}
\def\fm@cursnom{}

% clau del curs, color i color del text damunt del color
\newcommand{\fm@curs}[3]{\@namedef{fm@color@#1}{#2}\@namedef{fm@text@#1}{#3}}
\fm@curs{1eso}{E6B84A}{navy}
\fm@curs{2eso}{3B82F6}{white}
\fm@curs{3eso}{FF7A50}{white}
\fm@curs{4eso-a}{A78BFA}{white}
\fm@curs{4eso-b}{6D28D9}{white}
\fm@curs{1batx-mi}{9B1C1C}{white}
\fm@curs{1batx-ccss}{EF6B6B}{white}
\fm@curs{2batx-mii}{1E6B47}{white}
\fm@curs{2batx-ccss}{2BB3A3}{white}

\newcommand{\fmcurs}[1]{%
  \@ifundefined{fm@color@#1}{\PackageError{fentmates}{Curs desconegut: #1}{}}{}%
  \edef\fm@tmp{\noexpand\definecolor{curs}{HTML}{\@nameuse{fm@color@#1}}%
               \noexpand\colorlet{curstext}{\@nameuse{fm@text@#1}}}\fm@tmp
  \colorlet{cursfosc}{curs!72!black}%
  \def\fm@logo{\fm@logodir/vertical-#1.pdf}%
  \def\fm@marca{\fm@logodir/marca-#1.pdf}}

% ---------- capçalera de la primera pàgina ----------
% #1 matèria (petit), #2 curs, #3 tipus de document i unitat, #4 títol
\newcommand{\capcalera}[4]{%
  \gdef\fm@cursnom{#2}%
  \par\noindent
  \begin{tcolorbox}[colback=curs,colframe=curs,boxrule=0pt,arc=2.2mm,width=3.5cm,height=2.15cm,
      valign=center,left=3mm,right=2mm,top=1mm,bottom=1mm,nobeforeafter,box align=center]
    \color{curstext}{\footnotesize #1}\par\vspace{1pt}{\fontsize{19}{21}\selectfont\bfseries #2\par}%
  \end{tcolorbox}\hspace{5mm}%
  \begin{minipage}[c]{\dimexpr\linewidth-3.5cm-5mm-2.9cm\relax}
    {\color{cursfosc}\small\bfseries\MakeUppercase{#3}}\par\vspace{2pt}
    {\color{navy}\fontsize{20}{23}\selectfont\bfseries #4\par}
  \end{minipage}\hfill
  \begin{minipage}[c]{2.6cm}\raggedleft\includegraphics[width=2.6cm]{\fm@logo}\end{minipage}\par
  \vspace{2pt}}

% ---------- barres de secció del color del curs ----------
\renewcommand{\seccio}[2]{\par\needspace{\fmneed\baselineskip}\bigskip
  \begin{tcolorbox}[colback=curs,colframe=curs,coltext=curstext,arc=2mm,boxrule=0pt,left=3mm,right=3mm,top=1.2mm,bottom=1.2mm]
  {\large\bfseries #1\quad #2}
  \end{tcolorbox}\vspace{-2pt}}

% ---------- capçalera i peu de les pàgines ----------
% primera pàgina: logos institucionals al peu
\fancypagestyle{primera}{\fancyhf{}\renewcommand{\headrulewidth}{0pt}
  \fancyfoot[L]{\raisebox{-0.32\height}{\includegraphics[height=0.78cm]{logo-gva.png}}\hspace{5mm}%
    \raisebox{-0.32\height}{\includegraphics[height=0.78cm]{logo-vermellar.png}}}
  \fancyfoot[R]{\scriptsize\color{muted}\href{https://creativecommons.org/licenses/by-nc-sa/4.0/deed.ca}{\ccbyncsa}\ \ CC BY-NC-SA 4.0 · fentmates.org\qquad\small\thepage}}
% la resta: el símbol a l'esquerra i el curs a la dreta, amb el títol al centre com abans
\fancyhead[L]{\raisebox{-0.3\height}{\includegraphics[height=0.62cm]{\fm@marca}}}
\fancyhead[R]{\small\bfseries\color{cursfosc}\fm@cursnom}
\renewcommand{\headrule}{\color{curs!45}\hrule height \headrulewidth}
\makeatother

\fmcurs{3eso}

\begin{document}
\thispagestyle{primera}

\capcalera{Matemàtiques}{3r ESO}{Fitxa d'exercicis · Unitat 11}{Cossos geomètrics}
\vspace{10pt}
Nom i cognoms: \rule{8.4cm}{0.4pt}\hfill Grup: \rule{1.6cm}{0.4pt}\hfill Data: \rule{2.2cm}{0.4pt}
\vspace{8pt}

\begin{tcolorbox}[colback=soft,colframe=soft,arc=2mm,left=3mm,right=3mm,top=1.5mm,bottom=1.5mm]
\small Exercicis graduats per nivells de dificultat:\quad \estrelles{1}\ bàsic \qquad \estrelles{2}\ mitjà \qquad \estrelles{3}\ avançat.\\[2pt] Resol-los al quadern mostrant tots els passos i fes un dibuix de cada cos. Si un resultat no és exacte, arrodoneix-lo a les centèsimes; per a $\pi$, fes servir la tecla de la calculadora.
\end{tcolorbox}

% =====================================================
\seccio{1}{Poliedres. Fórmula d'Euler. Poliedres regulars}

\ex{1} Copia i completa la taula. Comprova en cada cas la fórmula d'Euler.\par\smallskip\begin{center}\renewcommand{\arraystretch}{1.25}\begin{tabular}{|l|C{1.6cm}|C{1.6cm}|C{1.6cm}|}\hline Poliedre & Cares & Vèrtexs & Arestes \\\hline prisma triangular & & &  \\\hline piràmide pentagonal & & &  \\\hline prisma hexagonal & & &  \\\hline piràmide octogonal & & &  \\\hline cub & & &  \\\hline\end{tabular}\end{center}\recorda{En un poliedre convex, \textbf{cares} + \textbf{vèrtexs} = \textbf{arestes} + 2: \ $C+V=A+2$ (fórmula d'Euler). Un prisma de base de $n$ costats té $n+2$ cares, $2n$ vèrtexs i $3n$ arestes; una piràmide, $n+1$ cares, $n+1$ vèrtexs i $2n$ arestes.}

\ex{1} Fes servir la fórmula d'Euler per a calcular el que falta en cada poliedre convex:\begin{tasks}(3)
  \task $C=8$, \ $V=12$, \ $A=\ ?$
  \task $V=10$, \ $A=15$, \ $C=\ ?$
  \task $C=12$, \ $A=30$, \ $V=\ ?$
\end{tasks}

\ex{1} Quins d'aquests poliedres són regulars? Justifica-ho.\begin{tasks}(3)
  \task tetraedre
  \task ortoedre
  \task octaedre
  \task piràmide quadrangular regular
  \task icosaedre
  \task cub
\end{tasks}

\ex{2} Un prisma té 24 arestes. Quin polígon té per base? Quantes cares i quants vèrtexs té?

\ex{2} Una piràmide té 11 vèrtexs. Quin polígon té per base? Quantes cares i arestes té?

\ex{2} Pot existir un poliedre convex amb: \ a) 7 cares, 7 vèrtexs i 12 arestes? \ b) 5 cares, 6 vèrtexs i 10 arestes?

\ex{3} Si ajuntem dues piràmides quadrangulars iguals per la base obtenim una \emph{bipiràmide}. Quantes cares, vèrtexs i arestes té? I si les piràmides són pentagonals?

\ex{3} Explica per què no hi pot haver cap poliedre regular amb cares hexagonals.


% =====================================================
\seccio{2}{Prismes: elements i àrea}

\ex{1} \begin{minipage}[t]{0.52\linewidth}Calcula l'àrea total:\begin{tasks}(1)
  \task d'un cub de 4 cm d'aresta;
  \task de l'ortoedre de la figura.
\end{tasks}\end{minipage}\hfill\begin{minipage}[t]{0.34\linewidth}\vspace{0pt}\centering \begin{tikzpicture}[x=0.45cm,y=0.45cm,baseline=(current bounding box.center),line join=round]
\path[fill=teal!14,draw=none] (0.000,3.709) -- (0.000,0.000) -- (5.298,-1.055) -- (5.298,2.654) -- cycle;
\path[fill=teal!14,draw=none] (5.298,2.654) -- (5.298,-1.055) -- (6.706,-0.063) -- (6.706,3.646) -- cycle;
\path[fill=teal!14,draw=none] (1.408,4.701) -- (0.000,3.709) -- (5.298,2.654) -- (6.706,3.646) -- cycle;
\draw[dashed,draw=muted] (0.000,0.000) -- (1.408,0.992);
\draw[dashed,draw=muted] (1.408,0.992) -- (6.706,-0.063);
\draw[dashed,draw=muted] (1.408,0.992) -- (1.408,4.701);
\draw[thick] (0.000,0.000) -- (5.298,-1.055);
\draw[thick] (5.298,-1.055) -- (6.706,-0.063);
\draw[thick] (0.000,0.000) -- (0.000,3.709);
\draw[thick] (5.298,-1.055) -- (5.298,2.654);
\draw[thick] (0.000,3.709) -- (5.298,2.654);
\draw[thick] (0.000,3.709) -- (1.408,4.701);
\draw[thick] (6.706,-0.063) -- (6.706,3.646);
\draw[thick] (5.298,2.654) -- (6.706,3.646);
\draw[thick] (1.408,4.701) -- (6.706,3.646);
\node[font=\small,inner sep=2pt,below] at (2.649,-0.528) {$6$\,cm};
\node[font=\small,inner sep=2pt,right] at (6.002,-0.559) {$3$\,cm};
\node[font=\small,inner sep=2pt,right] at (5.298,0.799) {$4$\,cm};
\end{tikzpicture}\end{minipage}\par\recorda{Àrea d'un prisma recte: $A_L=P_{\text{base}}\cdot h$, \ $A_T=A_L+2\,A_{\text{base}}$. Ortoedre d'arestes $a,\,b,\,c$: $A_T=2(ab+ac+bc)$ i la diagonal és $D=\sqrt{a^2+b^2+c^2}$.}

\ex{1} Calcula la diagonal d'un ortoedre de 12 cm × 4 cm × 3 cm.

\ex{1} Un prisma quadrangular regular té 5 cm d'aresta de la base i 12 cm d'altura. Calcula'n l'àrea lateral i l'àrea total.

\ex{2} Calcula l'àrea total d'un prisma triangular regular de 4 cm d'aresta de la base i 10 cm d'altura.

\ex{2} Calcula l'àrea total d'un prisma hexagonal regular de 6 cm d'aresta de la base i 8 cm d'altura.

\ex{2} Volem fer una caixa sense tapa de 40 cm × 30 cm de base i 20 cm d'altura. Quant cartó necessitem?

\ex{3} La diagonal d'un cub mesura $\sqrt{48}$ cm. Quant mesura l'aresta? Calcula'n l'àrea total.

\ex{3} La base d'un prisma recte és un triangle rectangle de catets 6 cm i 8 cm, i l'altura del prisma és 15 cm. Calcula'n l'àrea total.


% =====================================================
\seccio{3}{Piràmides i troncs de piràmide}

\ex{1} \begin{minipage}[t]{0.56\linewidth}La piràmide quadrangular regular de la figura té 10 cm d'aresta de la base i 12 cm d'altura. Calcula'n l'apotema i l'àrea total.\end{minipage}\hfill\begin{minipage}[t]{0.3\linewidth}\vspace{0pt}\centering \begin{tikzpicture}[x=0.5cm,y=0.5cm,baseline=(current bounding box.center),line join=round]
\path[fill=teal!14,draw=none] (0.000,3.338) -- (-2.029,-0.232) -- (0.620,-0.760) -- cycle;
\path[fill=teal!14,draw=none] (0.000,3.338) -- (0.620,-0.760) -- (2.029,0.232) -- cycle;
\draw[dashed,draw=muted] (-2.029,-0.232) -- (-0.620,0.760);
\draw[dashed,draw=muted] (2.029,0.232) -- (-0.620,0.760);
\draw[dashed,draw=muted] (-0.620,0.760) -- (0.000,3.338);
\draw[thick] (-2.029,-0.232) -- (0.620,-0.760);
\draw[thick] (0.620,-0.760) -- (2.029,0.232);
\draw[thick] (-2.029,-0.232) -- (0.000,3.338);
\draw[thick] (0.620,-0.760) -- (0.000,3.338);
\draw[thick] (2.029,0.232) -- (0.000,3.338);
\node[font=\small,inner sep=2pt,below] at (-0.704,-0.496) {$10$\,cm};
\draw[dashed,draw=accent,thick] (0.000,0.000) -- (0.000,3.338);
\draw[dashed,draw=accent,thick] (0.000,0.000) -- (1.324,-0.264);
\draw[draw=navy,thick] (1.324,-0.264) -- (0.000,3.338);
\fill (0.000,0.000) circle (1.3pt);
\node[font=\small,inner sep=2pt,left,text=accent] at (0.000,1.669) {$12$\,cm};
\end{tikzpicture}\end{minipage}\par\recorda{En una piràmide regular, l'altura $h$, l'apotema de la base $a_b$ i l'apotema de la piràmide $a_p$ formen un triangle rectangle: $a_p^2=h^2+a_b^2$. \ $A_L=\frac{P_{\text{base}}\cdot a_p}{2}$, \ $A_T=A_L+A_{\text{base}}$.}

\ex{1} Una piràmide hexagonal regular té 4 cm d'aresta de la base i 9 cm d'apotema. Calcula'n l'àrea lateral.

\ex{2} Una piràmide quadrangular regular té 6 cm d'aresta de la base i 5 cm d'apotema. Quina altura té? Calcula'n l'àrea total.

\ex{2} Una piràmide quadrangular regular té 8 cm d'aresta de la base i 6 cm d'altura. Calcula l'apotema, l'aresta lateral i l'àrea total.

\ex{2} Calcula l'àrea total d'un tetraedre regular de 6 cm d'aresta.

\ex{3} Un tronc de piràmide quadrangular regular té les arestes de les bases de 12 cm i 6 cm i una altura de 4 cm. Calcula l'apotema del tronc i l'àrea total.

\ex{3} La gran piràmide de Khufu (Gizeh) té una base quadrada de 230 m de costat i fa 146 m d'altura. Calcula l'apotema i l'àrea de les quatre cares laterals.


% =====================================================
\seccio{4}{Cilindre i con}

\ex{1} \begin{minipage}[t]{0.61\linewidth}Calcula l'àrea lateral i l'àrea total del cilindre de la figura.\end{minipage}\hfill\begin{minipage}[t]{0.25\linewidth}\vspace{0pt}\centering \begin{tikzpicture}[x=0.5cm,y=0.5cm,baseline=(current bounding box.center),line join=round]
\path[fill=teal!14,draw=none] (-1.500,2.967) -- (-1.500,0.000) -- (-1.497,-0.037) -- (-1.487,-0.073) -- (-1.471,-0.110) -- (-1.449,-0.145) -- (-1.420,-0.181) -- (-1.386,-0.215) -- (-1.345,-0.249) -- (-1.299,-0.281) -- (-1.247,-0.312) -- (-1.190,-0.342) -- (-1.128,-0.370) -- (-1.061,-0.397) -- (-0.989,-0.422) -- (-0.913,-0.446) -- (-0.833,-0.467) -- (-0.750,-0.487) -- (-0.663,-0.504) -- (-0.574,-0.519) -- (-0.482,-0.532) -- (-0.388,-0.543) -- (-0.293,-0.551) -- (-0.196,-0.557) -- (-0.098,-0.561) -- (-0.000,-0.562) -- (0.098,-0.561) -- (0.196,-0.557) -- (0.293,-0.551) -- (0.388,-0.543) -- (0.482,-0.532) -- (0.574,-0.519) -- (0.663,-0.504) -- (0.750,-0.487) -- (0.833,-0.467) -- (0.913,-0.446) -- (0.989,-0.422) -- (1.061,-0.397) -- (1.128,-0.370) -- (1.190,-0.342) -- (1.247,-0.312) -- (1.299,-0.281) -- (1.345,-0.249) -- (1.386,-0.215) -- (1.420,-0.181) -- (1.449,-0.145) -- (1.471,-0.110) -- (1.487,-0.073) -- (1.497,-0.037) -- (1.500,-0.000) -- (1.500,2.967) -- (1.500,2.967) -- (1.497,2.930) -- (1.487,2.894) -- (1.471,2.857) -- (1.449,2.822) -- (1.420,2.786) -- (1.386,2.752) -- (1.345,2.718) -- (1.299,2.686) -- (1.247,2.655) -- (1.190,2.625) -- (1.128,2.596) -- (1.061,2.570) -- (0.989,2.545) -- (0.913,2.521) -- (0.833,2.500) -- (0.750,2.480) -- (0.663,2.463) -- (0.574,2.448) -- (0.482,2.435) -- (0.388,2.424) -- (0.293,2.416) -- (0.196,2.410) -- (0.098,2.406) -- (-0.000,2.405) -- (-0.098,2.406) -- (-0.196,2.410) -- (-0.293,2.416) -- (-0.388,2.424) -- (-0.482,2.435) -- (-0.574,2.448) -- (-0.663,2.463) -- (-0.750,2.480) -- (-0.833,2.500) -- (-0.913,2.521) -- (-0.989,2.545) -- (-1.061,2.570) -- (-1.128,2.596) -- (-1.190,2.625) -- (-1.247,2.655) -- (-1.299,2.686) -- (-1.345,2.718) -- (-1.386,2.752) -- (-1.420,2.786) -- (-1.449,2.822) -- (-1.471,2.857) -- (-1.487,2.894) -- (-1.497,2.930) -- (-1.500,2.967) -- cycle;
\draw[dashed,draw=muted] (1.500,0.000) -- (1.497,0.037) -- (1.487,0.073) -- (1.471,0.110) -- (1.449,0.145) -- (1.420,0.181) -- (1.386,0.215) -- (1.345,0.249) -- (1.299,0.281) -- (1.247,0.312) -- (1.190,0.342) -- (1.128,0.370) -- (1.061,0.397) -- (0.989,0.422) -- (0.913,0.446) -- (0.833,0.467) -- (0.750,0.487) -- (0.663,0.504) -- (0.574,0.519) -- (0.482,0.532) -- (0.388,0.543) -- (0.293,0.551) -- (0.196,0.557) -- (0.098,0.561) -- (0.000,0.562) -- (-0.098,0.561) -- (-0.196,0.557) -- (-0.293,0.551) -- (-0.388,0.543) -- (-0.482,0.532) -- (-0.574,0.519) -- (-0.663,0.504) -- (-0.750,0.487) -- (-0.833,0.467) -- (-0.913,0.446) -- (-0.989,0.422) -- (-1.061,0.397) -- (-1.128,0.370) -- (-1.190,0.342) -- (-1.247,0.312) -- (-1.299,0.281) -- (-1.345,0.249) -- (-1.386,0.215) -- (-1.420,0.181) -- (-1.449,0.145) -- (-1.471,0.110) -- (-1.487,0.073) -- (-1.497,0.037) -- (-1.500,0.000);
\draw[thick] (-1.500,0.000) -- (-1.497,-0.037) -- (-1.487,-0.073) -- (-1.471,-0.110) -- (-1.449,-0.145) -- (-1.420,-0.181) -- (-1.386,-0.215) -- (-1.345,-0.249) -- (-1.299,-0.281) -- (-1.247,-0.312) -- (-1.190,-0.342) -- (-1.128,-0.370) -- (-1.061,-0.397) -- (-0.989,-0.422) -- (-0.913,-0.446) -- (-0.833,-0.467) -- (-0.750,-0.487) -- (-0.663,-0.504) -- (-0.574,-0.519) -- (-0.482,-0.532) -- (-0.388,-0.543) -- (-0.293,-0.551) -- (-0.196,-0.557) -- (-0.098,-0.561) -- (-0.000,-0.562) -- (0.098,-0.561) -- (0.196,-0.557) -- (0.293,-0.551) -- (0.388,-0.543) -- (0.482,-0.532) -- (0.574,-0.519) -- (0.663,-0.504) -- (0.750,-0.487) -- (0.833,-0.467) -- (0.913,-0.446) -- (0.989,-0.422) -- (1.061,-0.397) -- (1.128,-0.370) -- (1.190,-0.342) -- (1.247,-0.312) -- (1.299,-0.281) -- (1.345,-0.249) -- (1.386,-0.215) -- (1.420,-0.181) -- (1.449,-0.145) -- (1.471,-0.110) -- (1.487,-0.073) -- (1.497,-0.037) -- (1.500,-0.000);
\draw[thick] (-1.500,0.000) -- (-1.500,2.967);
\draw[thick] (1.500,0.000) -- (1.500,2.967);
\draw[fill=teal!26,draw=black,thick] (1.500,2.967) -- (1.487,3.040) -- (1.449,3.112) -- (1.386,3.182) -- (1.299,3.248) -- (1.190,3.309) -- (1.061,3.364) -- (0.913,3.413) -- (0.750,3.454) -- (0.574,3.486) -- (0.388,3.510) -- (0.196,3.524) -- (0.000,3.529) -- (-0.196,3.524) -- (-0.388,3.510) -- (-0.574,3.486) -- (-0.750,3.454) -- (-0.913,3.413) -- (-1.061,3.364) -- (-1.190,3.309) -- (-1.299,3.248) -- (-1.386,3.182) -- (-1.449,3.112) -- (-1.487,3.040) -- (-1.500,2.967) -- (-1.487,2.894) -- (-1.449,2.822) -- (-1.386,2.752) -- (-1.299,2.686) -- (-1.190,2.625) -- (-1.061,2.570) -- (-0.913,2.521) -- (-0.750,2.480) -- (-0.574,2.448) -- (-0.388,2.424) -- (-0.196,2.410) -- (-0.000,2.405) -- (0.196,2.410) -- (0.388,2.424) -- (0.574,2.448) -- (0.750,2.480) -- (0.913,2.521) -- (1.061,2.570) -- (1.190,2.625) -- (1.299,2.686) -- (1.386,2.752) -- (1.449,2.822) -- (1.487,2.894) -- (1.500,2.967) -- cycle;
\draw[dashed,draw=accent,thick] (0.000,0.000) -- (1.500,0.000);
\fill (0.000,0.000) circle (1.3pt);
\node[font=\small,inner sep=2pt,above,text=accent] at (0.750,0.000) {$3$\,cm};
\node[font=\small,inner sep=2pt,right] at (1.500,1.483) {$10$\,cm};
\end{tikzpicture}\end{minipage}\par\recorda{Cilindre: $A_L=2\pi rh$, \ $A_T=2\pi rh+2\pi r^2$. \ Con: $g^2=h^2+r^2$, \ $A_L=\pi rg$, \ $A_T=\pi rg+\pi r^2$. Fes servir la tecla $\pi$ i arrodoneix a les centèsimes.}

\ex{1} \begin{minipage}[t]{0.61\linewidth}Calcula la generatriu, l'àrea lateral i l'àrea total del con de la figura.\end{minipage}\hfill\begin{minipage}[t]{0.25\linewidth}\vspace{0pt}\centering \begin{tikzpicture}[x=0.5cm,y=0.5cm,baseline=(current bounding box.center),line join=round]
\path[fill=teal!14,draw=none] (0.000,3.152) -- (-1.571,0.114) -- (-1.589,0.070) -- (-1.598,0.027) -- (-1.599,-0.017) -- (-1.592,-0.061) -- (-1.575,-0.105) -- (-1.551,-0.148) -- (-1.517,-0.190) -- (-1.476,-0.231) -- (-1.427,-0.271) -- (-1.370,-0.310) -- (-1.306,-0.346) -- (-1.234,-0.381) -- (-1.156,-0.414) -- (-1.072,-0.445) -- (-0.982,-0.473) -- (-0.887,-0.499) -- (-0.787,-0.522) -- (-0.682,-0.542) -- (-0.574,-0.559) -- (-0.463,-0.574) -- (-0.350,-0.585) -- (-0.234,-0.593) -- (-0.117,-0.598) -- (-0.000,-0.599) -- (0.117,-0.598) -- (0.234,-0.593) -- (0.350,-0.585) -- (0.463,-0.574) -- (0.574,-0.559) -- (0.682,-0.542) -- (0.787,-0.522) -- (0.887,-0.499) -- (0.982,-0.473) -- (1.072,-0.445) -- (1.156,-0.414) -- (1.234,-0.381) -- (1.306,-0.346) -- (1.370,-0.310) -- (1.427,-0.271) -- (1.476,-0.231) -- (1.517,-0.190) -- (1.551,-0.148) -- (1.575,-0.105) -- (1.592,-0.061) -- (1.599,-0.017) -- (1.598,0.027) -- (1.589,0.070) -- (1.571,0.114) -- cycle;
\draw[dashed,draw=muted] (1.571,0.114) -- (1.551,0.148) -- (1.526,0.181) -- (1.495,0.213) -- (1.460,0.245) -- (1.420,0.276) -- (1.375,0.306) -- (1.326,0.335) -- (1.273,0.363) -- (1.215,0.390) -- (1.153,0.416) -- (1.087,0.440) -- (1.018,0.462) -- (0.946,0.484) -- (0.870,0.503) -- (0.791,0.521) -- (0.710,0.537) -- (0.627,0.552) -- (0.541,0.564) -- (0.454,0.575) -- (0.365,0.584) -- (0.275,0.590) -- (0.184,0.595) -- (0.092,0.598) -- (0.000,0.599) -- (-0.092,0.598) -- (-0.184,0.595) -- (-0.275,0.590) -- (-0.365,0.584) -- (-0.454,0.575) -- (-0.541,0.564) -- (-0.627,0.552) -- (-0.710,0.537) -- (-0.791,0.521) -- (-0.870,0.503) -- (-0.946,0.484) -- (-1.018,0.462) -- (-1.087,0.440) -- (-1.153,0.416) -- (-1.215,0.390) -- (-1.273,0.363) -- (-1.326,0.335) -- (-1.375,0.306) -- (-1.420,0.276) -- (-1.460,0.245) -- (-1.495,0.213) -- (-1.526,0.181) -- (-1.551,0.148) -- (-1.571,0.114);
\draw[thick] (-1.571,0.114) -- (-1.589,0.070) -- (-1.598,0.027) -- (-1.599,-0.017) -- (-1.592,-0.061) -- (-1.575,-0.105) -- (-1.551,-0.148) -- (-1.517,-0.190) -- (-1.476,-0.231) -- (-1.427,-0.271) -- (-1.370,-0.310) -- (-1.306,-0.346) -- (-1.234,-0.381) -- (-1.156,-0.414) -- (-1.072,-0.445) -- (-0.982,-0.473) -- (-0.887,-0.499) -- (-0.787,-0.522) -- (-0.682,-0.542) -- (-0.574,-0.559) -- (-0.463,-0.574) -- (-0.350,-0.585) -- (-0.234,-0.593) -- (-0.117,-0.598) -- (-0.000,-0.599) -- (0.117,-0.598) -- (0.234,-0.593) -- (0.350,-0.585) -- (0.463,-0.574) -- (0.574,-0.559) -- (0.682,-0.542) -- (0.787,-0.522) -- (0.887,-0.499) -- (0.982,-0.473) -- (1.072,-0.445) -- (1.156,-0.414) -- (1.234,-0.381) -- (1.306,-0.346) -- (1.370,-0.310) -- (1.427,-0.271) -- (1.476,-0.231) -- (1.517,-0.190) -- (1.551,-0.148) -- (1.575,-0.105) -- (1.592,-0.061) -- (1.599,-0.017) -- (1.598,0.027) -- (1.589,0.070) -- (1.571,0.114);
\draw[thick] (0.000,3.152) -- (-1.571,0.114);
\draw[thick] (0.000,3.152) -- (1.571,0.114);
\draw[dashed,draw=accent,thick] (0.000,0.000) -- (1.600,0.000);
\draw[dashed,draw=accent,thick] (0.000,0.000) -- (0.000,3.152);
\fill (0.000,0.000) circle (1.3pt);
\node[font=\small,inner sep=2pt,above,text=accent] at (0.800,0.000) {$5$\,cm};
\node[font=\small,inner sep=2pt,left,text=accent] at (0.000,1.576) {$12$\,cm};
\node[font=\small,inner sep=2pt,right] at (0.800,1.576) {$g$};
\end{tikzpicture}\end{minipage}\par

\ex{1} Un con té 6 cm de radi i 10 cm de generatriu. Quina altura té? Calcula'n l'àrea total.

\ex{2} Una llanda de refresc té 8 cm de diàmetre i 12 cm d'altura. Quanta xapa cal per a fer-la?

\ex{2} Volem fer un barret de pallasso en forma de con de 9 cm de radi i 12 cm d'altura. Quant paper necessitem?

\ex{2} Un cilindre té una àrea lateral de $120\pi$ cm² i 10 cm d'altura. Quin radi té? Calcula'n l'àrea total.

\ex{3} Un con té 4 cm de radi i 12 cm de generatriu. Quin angle té el sector circular del seu desenvolupament?

\ex{3} Cal pintar la superfície lateral de 12 columnes cilíndriques de 0,4 m de radi i 5 m d'altura. Si un litre de pintura cobreix 4 m², quants litres calen?


% =====================================================
\seccio{5}{L'esfera. La Terra}

\ex{1} \begin{minipage}[t]{0.64\linewidth}Calcula l'àrea de:\begin{tasks}(1)
  \task una esfera de 4 cm de radi;
  \task una pilota de 10 cm de diàmetre;
  \task una bola de 0,5 m de radi.
\end{tasks}\end{minipage}\hfill\begin{minipage}[t]{0.22\linewidth}\vspace{0pt}\centering \begin{tikzpicture}[x=0.5cm,y=0.5cm,baseline=(current bounding box.center),line join=round]
\draw[fill=teal!14,draw=black,thick] (0.000,0.000) circle[radius=1.6];
\draw[dashed,draw=muted] (1.600,0.000) -- (1.597,0.039) -- (1.586,0.078) -- (1.569,0.117) -- (1.545,0.155) -- (1.515,0.193) -- (1.478,0.229) -- (1.435,0.265) -- (1.386,0.300) -- (1.330,0.333) -- (1.269,0.365) -- (1.203,0.395) -- (1.131,0.424) -- (1.055,0.451) -- (0.974,0.476) -- (0.889,0.498) -- (0.800,0.519) -- (0.708,0.538) -- (0.612,0.554) -- (0.514,0.568) -- (0.414,0.579) -- (0.312,0.588) -- (0.209,0.594) -- (0.105,0.598) -- (0.000,0.599) -- (-0.105,0.598) -- (-0.209,0.594) -- (-0.312,0.588) -- (-0.414,0.579) -- (-0.514,0.568) -- (-0.612,0.554) -- (-0.708,0.538) -- (-0.800,0.519) -- (-0.889,0.498) -- (-0.974,0.476) -- (-1.055,0.451) -- (-1.131,0.424) -- (-1.203,0.395) -- (-1.269,0.365) -- (-1.330,0.333) -- (-1.386,0.300) -- (-1.435,0.265) -- (-1.478,0.229) -- (-1.515,0.193) -- (-1.545,0.155) -- (-1.569,0.117) -- (-1.586,0.078) -- (-1.597,0.039) -- (-1.600,0.000);
\draw[thick] (-1.600,0.000) -- (-1.597,-0.039) -- (-1.586,-0.078) -- (-1.569,-0.117) -- (-1.545,-0.155) -- (-1.515,-0.193) -- (-1.478,-0.229) -- (-1.435,-0.265) -- (-1.386,-0.300) -- (-1.330,-0.333) -- (-1.269,-0.365) -- (-1.203,-0.395) -- (-1.131,-0.424) -- (-1.055,-0.451) -- (-0.974,-0.476) -- (-0.889,-0.498) -- (-0.800,-0.519) -- (-0.708,-0.538) -- (-0.612,-0.554) -- (-0.514,-0.568) -- (-0.414,-0.579) -- (-0.312,-0.588) -- (-0.209,-0.594) -- (-0.105,-0.598) -- (-0.000,-0.599) -- (0.105,-0.598) -- (0.209,-0.594) -- (0.312,-0.588) -- (0.414,-0.579) -- (0.514,-0.568) -- (0.612,-0.554) -- (0.708,-0.538) -- (0.800,-0.519) -- (0.889,-0.498) -- (0.974,-0.476) -- (1.055,-0.451) -- (1.131,-0.424) -- (1.203,-0.395) -- (1.269,-0.365) -- (1.330,-0.333) -- (1.386,-0.300) -- (1.435,-0.265) -- (1.478,-0.229) -- (1.515,-0.193) -- (1.545,-0.155) -- (1.569,-0.117) -- (1.586,-0.078) -- (1.597,-0.039) -- (1.600,-0.000);
\draw[dashed,draw=accent,thick] (0.000,0.000) -- (1.600,0.000);
\fill (0.000,0.000) circle (1.3pt);
\node[font=\small,inner sep=2pt,above,text=accent] at (0.800,0.000) {$r$};
\end{tikzpicture}\end{minipage}\par\recorda{Àrea de l'esfera: $A=4\pi r^2$ (quatre cercles màxims). \ La \textbf{latitud} va de $0^\circ$ a $90^\circ$ N o S (des de l'equador) i la \textbf{longitud}, de $0^\circ$ a $180^\circ$ E o O (des del meridià de Greenwich).}

\ex{2} Una esfera té una àrea de $144\pi$ cm². Quin radi té?

\ex{2} Una cúpula té forma de semiesfera de 5 m de radi. Quina superfície cal impermeabilitzar per fora?

\ex{2} Dues ciutats estan sobre el mateix meridià, l'una a $40^\circ$ N i l'altra a $10^\circ$ S. Quants graus de latitud les separen? Si un meridià sencer (la volta a la Terra) fa 40\,000 km, quina distància hi ha entre elles sobre el meridià?

\ex{2} Quan a Greenwich són les 12 hores solars, quina hora solar és en un lloc de longitud $45^\circ$ E? I en un de longitud $30^\circ$ O?

\ex{3} El radi de la Terra és 6371 km i el de la Lluna, 1737 km. Quantes vegades més gran és la superfície de la Terra que la de la Lluna?


% =====================================================
\seccio{6}{Volum. Prismes i cilindres}

\ex{1} Expressa en la unitat indicada:\begin{tasks}(3)
  \task 3,2 m³ en L
  \task 450 cm³ en L
  \task 2,5 L en cm³
  \task 0,75 dm³ en mL
  \task 12\,000 L en m³
  \task 5 hL en m³
\end{tasks}\recorda{Cada unitat de volum és 1000 vegades la següent: $1\text{ m}^3=1000\text{ dm}^3$, \ $1\text{ dm}^3=1000\text{ cm}^3$. \ $1\text{ dm}^3=1\text{ L}$, \ $1\text{ cm}^3=1\text{ mL}$. \ Prisma i cilindre: $V=A_{\text{base}}\cdot h$; \ $V_{\text{cilindre}}=\pi r^2h$.}

\ex{1} Calcula el volum de:\begin{tasks}(1)
  \task un cub de 6 cm d'aresta;
  \task un ortoedre de 8 cm × 5 cm × 3 cm;
  \task un cilindre de 4 cm de radi i 10 cm d'altura.
\end{tasks}

\ex{1} La base d'un prisma és un triangle rectangle de catets 5 cm i 12 cm. Si el prisma fa 20 cm d'altura, quin volum té?

\ex{2} Calcula el volum d'un prisma hexagonal regular de 2 cm d'aresta de la base i 5 cm d'altura.

\ex{2} Un aquari fa 60 cm × 30 cm de base i 40 cm d'altura. Quants litres hi caben? Quants litres hi ha si l'aigua arriba a $\frac34$ de l'altura?

\ex{2} Un dipòsit cilíndric té 1,5 m de radi i 4 m d'altura. Quants litres d'aigua hi caben?

\ex{2} Una piscina fa 25 m de llarg i 10 m d'ample. El fons és inclinat: la fondària va d'1 m a 3 m. Quin volum d'aigua hi cap? (És un prisma amb una base en forma de trapezi.)

\ex{3} Un got cilíndric fa 4 cm de radi. Si hi aboquem mig litre d'aigua, fins a quina altura arriba?

\ex{3} Una rajola fa 24 cm × 11,5 cm × 5 cm i té tres forats cilíndrics de 1,5 cm de radi que la travessen (en la direcció dels 5 cm). Quin volum de fang té?


% =====================================================
\seccio{7}{Volum de piràmides, cons i esferes}

\ex{1} Calcula el volum de:\begin{tasks}(1)
  \task una piràmide quadrangular de 6 cm d'aresta de la base i 10 cm d'altura;
  \task un con de 3 cm de radi i 7 cm d'altura;
  \task una esfera de 6 cm de radi.
\end{tasks}\recorda{$V_{\text{piràmide}}=\frac{A_{\text{base}}\cdot h}{3}$; \ $V_{\text{con}}=\frac{\pi r^2h}{3}$; \ $V_{\text{esfera}}=\frac43\pi r^3$.}

\ex{1} Quin volum té una bola de 20 cm de diàmetre? Expressa'l en cm³ i en litres.

\ex{2} Una piràmide quadrangular regular té 10 cm d'aresta de la base i 13 cm d'apotema. Calcula'n l'altura i el volum.

\ex{2} Un con té 6 cm de radi i 10 cm de generatriu. Calcula'n el volum.

\ex{2} Quants litres d'aire hi ha dins d'una pilota de bàsquet d'11 cm de radi?

\ex{2} Calcula el volum d'una piràmide hexagonal regular de 4 cm d'aresta de la base i 9 cm d'altura.

\ex{3} Un bol té forma de semiesfera de 12 cm de radi interior. Quants litres hi caben?

\ex{3} Una esfera i un con tenen el mateix radi, 6 cm, i el mateix volum. Quina altura té el con?


% =====================================================
\seccio{8}{Cossos compostos i problemes}

\ex{1} Una peça està formada per un cub de 6 cm d'aresta i, al damunt, una piràmide de la mateixa base i 4 cm d'altura. Calcula'n el volum.\recorda{Descompon el cos en cossos simples. El volum és la suma (o la diferència, si hi ha forats) dels volums. Per a l'àrea, suma només les superfícies que queden per fora.}

\ex{2} \begin{minipage}[t]{0.61\linewidth}Una sitja té la forma de la figura: un cilindre de 3 m de radi i 10 m d'altura amb una teulada cònica de 4 m d'altura. Quin volum té?\end{minipage}\hfill\begin{minipage}[t]{0.25\linewidth}\vspace{0pt}\centering \begin{tikzpicture}[x=0.42cm,y=0.42cm,baseline=(current bounding box.center),line join=round]
\path[fill=teal!14,draw=none] (-1.500,3.152) -- (-1.500,0.000) -- (-1.497,-0.037) -- (-1.487,-0.073) -- (-1.471,-0.110) -- (-1.449,-0.145) -- (-1.420,-0.181) -- (-1.386,-0.215) -- (-1.345,-0.249) -- (-1.299,-0.281) -- (-1.247,-0.312) -- (-1.190,-0.342) -- (-1.128,-0.370) -- (-1.061,-0.397) -- (-0.989,-0.422) -- (-0.913,-0.446) -- (-0.833,-0.467) -- (-0.750,-0.487) -- (-0.663,-0.504) -- (-0.574,-0.519) -- (-0.482,-0.532) -- (-0.388,-0.543) -- (-0.293,-0.551) -- (-0.196,-0.557) -- (-0.098,-0.561) -- (-0.000,-0.562) -- (0.098,-0.561) -- (0.196,-0.557) -- (0.293,-0.551) -- (0.388,-0.543) -- (0.482,-0.532) -- (0.574,-0.519) -- (0.663,-0.504) -- (0.750,-0.487) -- (0.833,-0.467) -- (0.913,-0.446) -- (0.989,-0.422) -- (1.061,-0.397) -- (1.128,-0.370) -- (1.190,-0.342) -- (1.247,-0.312) -- (1.299,-0.281) -- (1.345,-0.249) -- (1.386,-0.215) -- (1.420,-0.181) -- (1.449,-0.145) -- (1.471,-0.110) -- (1.487,-0.073) -- (1.497,-0.037) -- (1.500,-0.000) -- (1.500,3.152) -- (1.500,3.152) -- (1.497,3.116) -- (1.487,3.079) -- (1.471,3.043) -- (1.449,3.007) -- (1.420,2.972) -- (1.386,2.937) -- (1.345,2.904) -- (1.299,2.871) -- (1.247,2.840) -- (1.190,2.810) -- (1.128,2.782) -- (1.061,2.755) -- (0.989,2.730) -- (0.913,2.707) -- (0.833,2.685) -- (0.750,2.666) -- (0.663,2.648) -- (0.574,2.633) -- (0.482,2.620) -- (0.388,2.610) -- (0.293,2.601) -- (0.196,2.595) -- (0.098,2.592) -- (-0.000,2.591) -- (-0.098,2.592) -- (-0.196,2.595) -- (-0.293,2.601) -- (-0.388,2.610) -- (-0.482,2.620) -- (-0.574,2.633) -- (-0.663,2.648) -- (-0.750,2.666) -- (-0.833,2.685) -- (-0.913,2.707) -- (-0.989,2.730) -- (-1.061,2.755) -- (-1.128,2.782) -- (-1.190,2.810) -- (-1.247,2.840) -- (-1.299,2.871) -- (-1.345,2.904) -- (-1.386,2.937) -- (-1.420,2.972) -- (-1.449,3.007) -- (-1.471,3.043) -- (-1.487,3.079) -- (-1.497,3.116) -- (-1.500,3.152) -- cycle;
\draw[dashed,draw=muted] (1.500,0.000) -- (1.497,0.037) -- (1.487,0.073) -- (1.471,0.110) -- (1.449,0.145) -- (1.420,0.181) -- (1.386,0.215) -- (1.345,0.249) -- (1.299,0.281) -- (1.247,0.312) -- (1.190,0.342) -- (1.128,0.370) -- (1.061,0.397) -- (0.989,0.422) -- (0.913,0.446) -- (0.833,0.467) -- (0.750,0.487) -- (0.663,0.504) -- (0.574,0.519) -- (0.482,0.532) -- (0.388,0.543) -- (0.293,0.551) -- (0.196,0.557) -- (0.098,0.561) -- (0.000,0.562) -- (-0.098,0.561) -- (-0.196,0.557) -- (-0.293,0.551) -- (-0.388,0.543) -- (-0.482,0.532) -- (-0.574,0.519) -- (-0.663,0.504) -- (-0.750,0.487) -- (-0.833,0.467) -- (-0.913,0.446) -- (-0.989,0.422) -- (-1.061,0.397) -- (-1.128,0.370) -- (-1.190,0.342) -- (-1.247,0.312) -- (-1.299,0.281) -- (-1.345,0.249) -- (-1.386,0.215) -- (-1.420,0.181) -- (-1.449,0.145) -- (-1.471,0.110) -- (-1.487,0.073) -- (-1.497,0.037) -- (-1.500,0.000);
\draw[thick] (-1.500,0.000) -- (-1.497,-0.037) -- (-1.487,-0.073) -- (-1.471,-0.110) -- (-1.449,-0.145) -- (-1.420,-0.181) -- (-1.386,-0.215) -- (-1.345,-0.249) -- (-1.299,-0.281) -- (-1.247,-0.312) -- (-1.190,-0.342) -- (-1.128,-0.370) -- (-1.061,-0.397) -- (-0.989,-0.422) -- (-0.913,-0.446) -- (-0.833,-0.467) -- (-0.750,-0.487) -- (-0.663,-0.504) -- (-0.574,-0.519) -- (-0.482,-0.532) -- (-0.388,-0.543) -- (-0.293,-0.551) -- (-0.196,-0.557) -- (-0.098,-0.561) -- (-0.000,-0.562) -- (0.098,-0.561) -- (0.196,-0.557) -- (0.293,-0.551) -- (0.388,-0.543) -- (0.482,-0.532) -- (0.574,-0.519) -- (0.663,-0.504) -- (0.750,-0.487) -- (0.833,-0.467) -- (0.913,-0.446) -- (0.989,-0.422) -- (1.061,-0.397) -- (1.128,-0.370) -- (1.190,-0.342) -- (1.247,-0.312) -- (1.299,-0.281) -- (1.345,-0.249) -- (1.386,-0.215) -- (1.420,-0.181) -- (1.449,-0.145) -- (1.471,-0.110) -- (1.487,-0.073) -- (1.497,-0.037) -- (1.500,-0.000);
\draw[thick] (-1.500,0.000) -- (-1.500,3.152);
\draw[thick] (1.500,0.000) -- (1.500,3.152);
\path[fill=teal!14,draw=none] (0.000,4.450) -- (-1.352,3.396) -- (-1.402,3.352) -- (-1.442,3.307) -- (-1.471,3.262) -- (-1.491,3.215) -- (-1.499,3.168) -- (-1.498,3.120) -- (-1.485,3.073) -- (-1.462,3.027) -- (-1.429,2.981) -- (-1.385,2.937) -- (-1.332,2.894) -- (-1.270,2.853) -- (-1.198,2.814) -- (-1.118,2.778) -- (-1.030,2.744) -- (-0.935,2.713) -- (-0.833,2.685) -- (-0.725,2.661) -- (-0.612,2.639) -- (-0.495,2.622) -- (-0.374,2.608) -- (-0.251,2.598) -- (-0.126,2.593) -- (-0.000,2.591) -- (0.126,2.593) -- (0.251,2.598) -- (0.374,2.608) -- (0.495,2.622) -- (0.612,2.639) -- (0.725,2.661) -- (0.833,2.685) -- (0.935,2.713) -- (1.030,2.744) -- (1.118,2.778) -- (1.198,2.814) -- (1.270,2.853) -- (1.332,2.894) -- (1.385,2.937) -- (1.429,2.981) -- (1.462,3.027) -- (1.485,3.073) -- (1.498,3.120) -- (1.499,3.168) -- (1.491,3.215) -- (1.471,3.262) -- (1.442,3.307) -- (1.402,3.352) -- (1.352,3.396) -- cycle;
\draw[dashed,draw=muted] (1.352,3.396) -- (1.320,3.419) -- (1.286,3.442) -- (1.248,3.464) -- (1.208,3.486) -- (1.165,3.506) -- (1.119,3.527) -- (1.071,3.546) -- (1.021,3.564) -- (0.969,3.581) -- (0.914,3.598) -- (0.857,3.614) -- (0.799,3.628) -- (0.738,3.642) -- (0.677,3.654) -- (0.613,3.665) -- (0.549,3.675) -- (0.483,3.684) -- (0.416,3.692) -- (0.348,3.699) -- (0.279,3.705) -- (0.210,3.709) -- (0.140,3.712) -- (0.070,3.714) -- (0.000,3.714) -- (-0.070,3.714) -- (-0.140,3.712) -- (-0.210,3.709) -- (-0.279,3.705) -- (-0.348,3.699) -- (-0.416,3.692) -- (-0.483,3.684) -- (-0.549,3.675) -- (-0.613,3.665) -- (-0.677,3.654) -- (-0.738,3.642) -- (-0.799,3.628) -- (-0.857,3.614) -- (-0.914,3.598) -- (-0.969,3.581) -- (-1.021,3.564) -- (-1.071,3.546) -- (-1.119,3.527) -- (-1.165,3.506) -- (-1.208,3.486) -- (-1.248,3.464) -- (-1.286,3.442) -- (-1.320,3.419) -- (-1.352,3.396);
\draw[thick] (-1.352,3.396) -- (-1.402,3.352) -- (-1.442,3.307) -- (-1.471,3.262) -- (-1.491,3.215) -- (-1.499,3.168) -- (-1.498,3.120) -- (-1.485,3.073) -- (-1.462,3.027) -- (-1.429,2.981) -- (-1.385,2.937) -- (-1.332,2.894) -- (-1.270,2.853) -- (-1.198,2.814) -- (-1.118,2.778) -- (-1.030,2.744) -- (-0.935,2.713) -- (-0.833,2.685) -- (-0.725,2.661) -- (-0.612,2.639) -- (-0.495,2.622) -- (-0.374,2.608) -- (-0.251,2.598) -- (-0.126,2.593) -- (-0.000,2.591) -- (0.126,2.593) -- (0.251,2.598) -- (0.374,2.608) -- (0.495,2.622) -- (0.612,2.639) -- (0.725,2.661) -- (0.833,2.685) -- (0.935,2.713) -- (1.030,2.744) -- (1.118,2.778) -- (1.198,2.814) -- (1.270,2.853) -- (1.332,2.894) -- (1.385,2.937) -- (1.429,2.981) -- (1.462,3.027) -- (1.485,3.073) -- (1.498,3.120) -- (1.499,3.168) -- (1.491,3.215) -- (1.471,3.262) -- (1.442,3.307) -- (1.402,3.352) -- (1.352,3.396);
\draw[thick] (0.000,4.450) -- (-1.352,3.396);
\draw[thick] (0.000,4.450) -- (1.352,3.396);
\draw[dashed,draw=accent,thick] (0.000,0.000) -- (1.500,0.000);
\node[font=\small,inner sep=2pt,below,text=accent] at (0.750,0.000) {$3$\,m};
\node[font=\small,inner sep=2pt,right] at (1.500,1.576) {$10$\,m};
\draw[dashed,draw=accent,thick] (0.000,3.152) -- (0.000,4.450);
\node[font=\small,inner sep=2pt,left,text=accent] at (-0.800,3.801) {$4$\,m};
\end{tikzpicture}\end{minipage}\par

\ex{2} \begin{minipage}[t]{0.66\linewidth}Una càpsula és un cilindre de 2 cm de radi i 6 cm de llarg amb una semiesfera a cada extrem. Calcula'n el volum i l'àrea.\end{minipage}\hfill\begin{minipage}[t]{0.2\linewidth}\vspace{0pt}\centering \begin{tikzpicture}[x=0.5cm,y=0.5cm,baseline=(current bounding box.center),line join=round]
\path[fill=teal!14,draw=none] (1.000,2.782) -- (0.998,2.847) -- (0.991,2.912) -- (0.981,2.977) -- (0.966,3.040) -- (0.947,3.103) -- (0.924,3.164) -- (0.897,3.224) -- (0.866,3.282) -- (0.831,3.337) -- (0.793,3.390) -- (0.752,3.441) -- (0.707,3.489) -- (0.659,3.533) -- (0.609,3.575) -- (0.556,3.613) -- (0.500,3.648) -- (0.442,3.678) -- (0.383,3.705) -- (0.321,3.728) -- (0.259,3.747) -- (0.195,3.762) -- (0.131,3.773) -- (0.065,3.779) -- (0.000,3.782) -- (-0.065,3.779) -- (-0.131,3.773) -- (-0.195,3.762) -- (-0.259,3.747) -- (-0.321,3.728) -- (-0.383,3.705) -- (-0.442,3.678) -- (-0.500,3.648) -- (-0.556,3.613) -- (-0.609,3.575) -- (-0.659,3.533) -- (-0.707,3.489) -- (-0.752,3.441) -- (-0.793,3.390) -- (-0.831,3.337) -- (-0.866,3.282) -- (-0.897,3.224) -- (-0.924,3.164) -- (-0.947,3.103) -- (-0.966,3.040) -- (-0.981,2.977) -- (-0.991,2.912) -- (-0.998,2.847) -- (-1.000,2.782) -- (-1.000,-0.000) -- (-0.998,-0.065) -- (-0.991,-0.131) -- (-0.981,-0.195) -- (-0.966,-0.259) -- (-0.947,-0.321) -- (-0.924,-0.383) -- (-0.897,-0.442) -- (-0.866,-0.500) -- (-0.831,-0.556) -- (-0.793,-0.609) -- (-0.752,-0.659) -- (-0.707,-0.707) -- (-0.659,-0.752) -- (-0.609,-0.793) -- (-0.556,-0.831) -- (-0.500,-0.866) -- (-0.442,-0.897) -- (-0.383,-0.924) -- (-0.321,-0.947) -- (-0.259,-0.966) -- (-0.195,-0.981) -- (-0.131,-0.991) -- (-0.065,-0.998) -- (-0.000,-1.000) -- (0.065,-0.998) -- (0.131,-0.991) -- (0.195,-0.981) -- (0.259,-0.966) -- (0.321,-0.947) -- (0.383,-0.924) -- (0.442,-0.897) -- (0.500,-0.866) -- (0.556,-0.831) -- (0.609,-0.793) -- (0.659,-0.752) -- (0.707,-0.707) -- (0.752,-0.659) -- (0.793,-0.609) -- (0.831,-0.556) -- (0.866,-0.500) -- (0.897,-0.442) -- (0.924,-0.383) -- (0.947,-0.321) -- (0.966,-0.259) -- (0.981,-0.195) -- (0.991,-0.131) -- (0.998,-0.065) -- (1.000,-0.000) -- cycle;
\draw[thick] (-1.000,0.000) -- (-0.998,-0.025) -- (-0.991,-0.049) -- (-0.981,-0.073) -- (-0.966,-0.097) -- (-0.947,-0.120) -- (-0.924,-0.143) -- (-0.897,-0.166) -- (-0.866,-0.187) -- (-0.831,-0.208) -- (-0.793,-0.228) -- (-0.752,-0.247) -- (-0.707,-0.265) -- (-0.659,-0.282) -- (-0.609,-0.297) -- (-0.556,-0.311) -- (-0.500,-0.324) -- (-0.442,-0.336) -- (-0.383,-0.346) -- (-0.321,-0.355) -- (-0.259,-0.362) -- (-0.195,-0.367) -- (-0.131,-0.371) -- (-0.065,-0.374) -- (-0.000,-0.375) -- (0.065,-0.374) -- (0.131,-0.371) -- (0.195,-0.367) -- (0.259,-0.362) -- (0.321,-0.355) -- (0.383,-0.346) -- (0.442,-0.336) -- (0.500,-0.324) -- (0.556,-0.311) -- (0.609,-0.297) -- (0.659,-0.282) -- (0.707,-0.265) -- (0.752,-0.247) -- (0.793,-0.228) -- (0.831,-0.208) -- (0.866,-0.187) -- (0.897,-0.166) -- (0.924,-0.143) -- (0.947,-0.120) -- (0.966,-0.097) -- (0.981,-0.073) -- (0.991,-0.049) -- (0.998,-0.025) -- (1.000,-0.000);
\draw[thick] (-1.000,2.782) -- (-0.998,2.757) -- (-0.991,2.733) -- (-0.981,2.708) -- (-0.966,2.685) -- (-0.947,2.661) -- (-0.924,2.638) -- (-0.897,2.616) -- (-0.866,2.594) -- (-0.831,2.573) -- (-0.793,2.554) -- (-0.752,2.535) -- (-0.707,2.517) -- (-0.659,2.500) -- (-0.609,2.484) -- (-0.556,2.470) -- (-0.500,2.457) -- (-0.442,2.446) -- (-0.383,2.435) -- (-0.321,2.427) -- (-0.259,2.420) -- (-0.195,2.414) -- (-0.131,2.410) -- (-0.065,2.408) -- (-0.000,2.407) -- (0.065,2.408) -- (0.131,2.410) -- (0.195,2.414) -- (0.259,2.420) -- (0.321,2.427) -- (0.383,2.435) -- (0.442,2.446) -- (0.500,2.457) -- (0.556,2.470) -- (0.609,2.484) -- (0.659,2.500) -- (0.707,2.517) -- (0.752,2.535) -- (0.793,2.554) -- (0.831,2.573) -- (0.866,2.594) -- (0.897,2.616) -- (0.924,2.638) -- (0.947,2.661) -- (0.966,2.685) -- (0.981,2.708) -- (0.991,2.733) -- (0.998,2.757) -- (1.000,2.782);
\draw[thick] (1.000,2.782) -- (0.998,2.847) -- (0.991,2.912) -- (0.981,2.977) -- (0.966,3.040) -- (0.947,3.103) -- (0.924,3.164) -- (0.897,3.224) -- (0.866,3.282) -- (0.831,3.337) -- (0.793,3.390) -- (0.752,3.441) -- (0.707,3.489) -- (0.659,3.533) -- (0.609,3.575) -- (0.556,3.613) -- (0.500,3.648) -- (0.442,3.678) -- (0.383,3.705) -- (0.321,3.728) -- (0.259,3.747) -- (0.195,3.762) -- (0.131,3.773) -- (0.065,3.779) -- (0.000,3.782) -- (-0.065,3.779) -- (-0.131,3.773) -- (-0.195,3.762) -- (-0.259,3.747) -- (-0.321,3.728) -- (-0.383,3.705) -- (-0.442,3.678) -- (-0.500,3.648) -- (-0.556,3.613) -- (-0.609,3.575) -- (-0.659,3.533) -- (-0.707,3.489) -- (-0.752,3.441) -- (-0.793,3.390) -- (-0.831,3.337) -- (-0.866,3.282) -- (-0.897,3.224) -- (-0.924,3.164) -- (-0.947,3.103) -- (-0.966,3.040) -- (-0.981,2.977) -- (-0.991,2.912) -- (-0.998,2.847) -- (-1.000,2.782) -- (-1.000,-0.000) -- (-0.998,-0.065) -- (-0.991,-0.131) -- (-0.981,-0.195) -- (-0.966,-0.259) -- (-0.947,-0.321) -- (-0.924,-0.383) -- (-0.897,-0.442) -- (-0.866,-0.500) -- (-0.831,-0.556) -- (-0.793,-0.609) -- (-0.752,-0.659) -- (-0.707,-0.707) -- (-0.659,-0.752) -- (-0.609,-0.793) -- (-0.556,-0.831) -- (-0.500,-0.866) -- (-0.442,-0.897) -- (-0.383,-0.924) -- (-0.321,-0.947) -- (-0.259,-0.966) -- (-0.195,-0.981) -- (-0.131,-0.991) -- (-0.065,-0.998) -- (-0.000,-1.000) -- (0.065,-0.998) -- (0.131,-0.991) -- (0.195,-0.981) -- (0.259,-0.966) -- (0.321,-0.947) -- (0.383,-0.924) -- (0.442,-0.897) -- (0.500,-0.866) -- (0.556,-0.831) -- (0.609,-0.793) -- (0.659,-0.752) -- (0.707,-0.707) -- (0.752,-0.659) -- (0.793,-0.609) -- (0.831,-0.556) -- (0.866,-0.500) -- (0.897,-0.442) -- (0.924,-0.383) -- (0.947,-0.321) -- (0.966,-0.259) -- (0.981,-0.195) -- (0.991,-0.131) -- (0.998,-0.065) -- (1.000,-0.000) -- (1.000,2.782);
\node[font=\small,inner sep=2pt,right] at (1.000,1.391) {$6$\,cm};
\draw[dashed,draw=accent,thick] (0.000,0.000) -- (1.000,0.000);
\node[font=\small,inner sep=2pt,below,text=accent] at (0.500,-1.000) {$2$\,cm};
\end{tikzpicture}\end{minipage}\par

\ex{2} Un gelat és un cucurutxo cònic de 2,5 cm de radi i 12 cm d'altura, ple, amb una semiesfera de gelat al damunt. Quin volum de gelat hi ha?

\ex{2} Un tub de 50 cm de llarg té 6 cm de radi exterior i 5 cm de radi interior. Quin volum de material té?

\ex{2} Un cub de 10 cm d'aresta té un forat cilíndric de 3 cm de radi que el travessa de banda a banda. Quin volum té?

\ex{3} Un cubell té forma de tronc de con: el fons fa 12 cm de radi, la boca 18 cm i l'altura és de 24 cm. Quants litres hi caben? \emph{(Pista: és un con gran al qual llevem un con menut; calcula l'altura del menut amb el teorema de Tales.)}

\ex{3} Deixem caure una bola d'acer de 3 cm de radi dins d'un got cilíndric de 6 cm de radi que conté aigua (i la bola queda coberta). Quant puja el nivell de l'aigua?

\ex{3} Un dipòsit d'aigua és un cilindre de 2 m de radi i 5 m d'altura amb el fons en forma de semiesfera. Quants m³ hi caben? Quant de temps tardarà a omplir-se amb una aixeta que aboca 50 L per minut?

\vfill
\begin{center}\color{muted}\small\faCheckCircle[regular]\ \ Ànim! Si et costa algun exercici, repassa els apunts interactius de la unitat.\end{center}

\end{document}
